\subsection{Earlier bounds}

Karatsuba and Ofman reduced a product to three products of half the
length. Toom and Cook split the inputs into more pieces and used
evaluation and interpolation to recover their product
\cite{karatsubaofman,toom,cook}. The nearly linear algorithms treat
multiplication as convolution. Sch\"onhage and Strassen obtained
$O(n\log n\log\log n)$ tape time using Fourier transforms over
residue rings \cite{schonhagestrassen}. F\"urer reduced the factor
above $n\log n$ to $2^{O(\log^*n)}$ \cite{furer}. Modular arithmetic
and short lattice vectors gave further constructions
\cite{DeKururSahaSaptharishi2013,HarveyHoevenLecerf2016,HarveyHoeven2019Lattice}.
Harvey and van der Hoeven reached $O(n\log n)$ on a multitape machine
using Gaussian resampling to obtain multidimensional Fourier transforms
of power-of-two side lengths \cite{hvh}.

An earlier fixed-tape algorithm combines resampling with a finite
network that shares butterfly operations and address rearrangements.
It gives worst-case time $O(n(1+\log n)^{1-2^{-182}})$
\cite[Theorem~1]{Tape}. Subsequent optimizations of the finite
schedules give larger numerical savings under explicit assumptions
on the cost of their component operations \cite{ChenTapeSchedules}.
Our implementation uses the original stream algorithms and address
changes. For the main convolution we use short cyclic factors
and apply the Fourier operations to many occupied rows at once.
The companion Fourier paper determines the exponent $1/2$ for exact
binary linear circuits, even when one circuit mixes the inputs of
independent transforms \cite[Theorem~\ext{F}{thm:main}]{oaF}. It
also describes earlier savings in operation counts
\cite{AlmanRao,FiniteFourier,ExplicitFourier,Sussman,Kartal}.
To obtain a tape algorithm we must move the rows into position,
compute the coefficients to sufficient precision, and recover every
product digit exactly.

For online multiplication, which must emit digits as the input arrives,
Paterson, Fischer and Meyer obtained an $\Omega(n\log n)$ bound even
on mean time, building on the overlap argument of Cook and Aanderaa
\cite{CA,PFM}. For moving data, Sto{\ss} proved an
$\Omega(N^2\log N)$ lower bound for transposing an $N\times N$ bit
matrix in a model that transports its entries intact \cite{moveonly}.
An ordinary tape machine may read both operands before writing the
product and may encode several input bits together before moving them.

Network coding gives a different route to lower bounds. Adler,
Harvey, Jain, Kleinberg and Lehman connected the undirected
multiple-unicast conjecture to matrix transposition
\cite{AHJ}. Assuming that conjecture, Afshani, Freksen, Kamma and
Larsen proved an $\Omega(n\log n)$ multiplication lower bound for
Boolean circuits with bounded fan-in and fan-out \cite{AFKL}.
Two independent constructions exhibit undirected networks where coding
exceeds the capacity of fractional routing \cite{ZhangLiLi,BravermanHe}. Our lower bound allows arbitrary local coding and charges for
communicating the resulting bits. It assumes variable-count Weak ETH,
which is equivalent to the Stearns--Hunt hypothesis for CNF satisfiability
\cite[Theorem~2.4]{WeakSH}. Subexponential algorithms for 3-SAT
\cite{Anthropic,falseeth} refute the original Exponential Time Hypothesis
\cite{IP} and are compatible with this weaker assumption.
